% Chapter4/chapter4.tex -- Implementation
% -----------------------------------------------------------------------------
% Chapter 4 takes a bottom-up approach: each section covers one computational
% primitive, discussing its implementation and pragma strategy before moving
% to the next. The top-level forward function is covered last, once all
% building blocks have been established.
%
% Section structure:
%   4.1  RMSNorm
%   4.2  Quantisation and Dequantisation
%   4.3  Matrix--Vector Multiplication
%   4.4  Softmax
%   4.5  Grouped-Query Attention
%   4.6  SwiGLU Feed-Forward Network
%   4.7  Top-Level Forward Function

\chapter{Implementation}
\label{ch:implementation}

All kernels in this chapter are implemented in Vitis HLS 2021.2,
targeting the \texttt{xczu9eg-ffvb1156-2-e} device on the ZCU102 evaluation
board at a design clock of \SI{200}{\mega\hertz}. Each primitive is first
presented as a standalone synthesisable function, allowing its resource
utilisation and initiation interval to be measured in isolation before it
is composed into the full transformer forward pass.

Development proceeded across two configurations. A scaled-down prototype
configuration was used to establish the kernel structure, verify the HLS
build flow, and confirm functional correctness in C-simulation without the
lengthy synthesis times and large on-chip resource requirements of the full
model. Once each component was validated at small scale, the configuration
constants were updated to target the full TinyLlama-1.1B dimensions.
Table~\ref{tab:config_comparison} summarises the two sets of parameters.

\begin{table}[H]
  \centering
  \caption{Prototype and full TinyLlama-1.1B configuration parameters
           used across kernel development.}
  \label{tab:config_comparison}
  \begin{tabular}{lrr}
    \toprule
    Parameter & Prototype & TinyLlama-1.1B \\
    \midrule
    Hidden dimension $d_{\text{model}}$ & 8       & 2{,}048  \\
    FFN hidden dimension $d_{\text{ff}}$ & 32     & 5{,}632  \\
    Number of layers $L$                 & 2       & 22       \\
    Query heads $H$                      & 2       & 32       \\
    KV heads $G$                         & 2       & 4        \\
    Vocabulary size $V$                  & 32      & 32{,}000 \\
    Maximum sequence length              & 16      & 2{,}048  \\
    Quantisation group size $GS$         & 4       & 64       \\
    \bottomrule
  \end{tabular}
\end{table}

The prototype dimensions were chosen deliberately to keep array sizes small
enough that HLS C-simulation completes in seconds and synthesis reports are
straightforward to inspect. With $d_{\text{model}} = 8$ and $GS = 4$, for
instance, every quantised vector contains exactly two groups, making
scale-factor indexing trivial to verify by hand. The constraint that $GS
\leq d_{\text{model}}$ must always hold, since quantising into groups larger
than the vector itself is undefined. At full scale the group size grows to
$GS = 64$, which divides evenly into all relevant dimensions in TinyLlama
(2048, 5632, 256) and aligns with a typical cache line size of 64 bytes,
making group-boundary loads naturally efficient.

% =============================================================================
\section{RMSNorm}
\label{sec:rmsnorm_impl}
% =============================================================================

\subsection{Implementation}
\label{subsec:rmsnorm_impl}

The \texttt{rmsnorm<S>} function is declared as a C++ template in
\texttt{forward.h}, parameterised by the vector length \texttt{S}. This
makes the same source compile correctly for both the prototype configuration
($S = d_{\text{model}} = 8$) and the full TinyLlama model
($S = d_{\text{model}} = 2048$) without any code changes. The function
takes three pointer arguments: an output array \texttt{o}, an input
activation array \texttt{x}, and a learned element-wise weight array
\texttt{weight}, all of length \texttt{S}.

The first design decision is to copy both \texttt{x} and \texttt{weight}
into local arrays \texttt{x\_buff} and \texttt{weight\_buff} before any
computation begins, as shown in Listing~\ref{lst:rmsnorm}. This is
necessary because Vitis HLS can only apply \texttt{ARRAY\_PARTITION} to
arrays whose storage it controls. Arrays received as pointer arguments may
alias external memory, which prevents the tool from freely splitting them
across multiple BRAM banks. By staging the data into local buffers first,
the \texttt{ARRAY\_PARTITION cyclic factor=4} pragma can be applied,
enabling four simultaneous reads per cycle in both computation loops.

\begin{lstlisting}[
  caption = {\texttt{rmsnorm<S>} implementation from \texttt{forward.h}.
             Data is staged into local buffers before computation so that
             \texttt{ARRAY\_PARTITION} can be applied. The reduction loop
             uses \texttt{UNROLL factor=4} to accumulate four partial sums
             simultaneously; the normalisation loop uses \texttt{UNROLL
             factor=64} since it is fully data-parallel.},
  label   = {lst:rmsnorm}
]
template <int S>
void rmsnorm(float o[S], float x[S], float weight[S])
{
    float x_buff[S], weight_buff[S], out_buff[S];
    #pragma HLS array_partition variable=x_buff      cyclic factor=4
    #pragma HLS array_partition variable=weight_buff cyclic factor=4
    #pragma HLS array_partition variable=out_buff    cyclic factor=4

    memcpy(x_buff,      x,      S * sizeof(float));
    memcpy(weight_buff, weight, S * sizeof(float));

    float ss = 0.0f;
sum_of_squares:
    for (int j = 0; j < S; j++) {
        #pragma HLS PIPELINE
        #pragma HLS UNROLL factor=4 skip_exit_check
        ss += x_buff[j] * x_buff[j];
    }
    ss /= S;
    ss += 1e-5f;
    ss = 1.0f / sqrtf(ss);   // reciprocal RMS

norm_and_scale:
    for (int j = 0; j < S; j++) {
        #pragma HLS PIPELINE
        #pragma HLS UNROLL factor=64
        out_buff[j] = weight_buff[j] * (ss * x_buff[j]);
    }
    memcpy(o, out_buff, S * sizeof(float));
}
\end{lstlisting}

The computation proceeds in two sequential loops with different parallelism
requirements. The first loop, \texttt{sum\_of\_squares}, accumulates the
sum of squared activations. This is a reduction with a loop-carried
dependency on the accumulator \texttt{ss}: each iteration reads the
current value of \texttt{ss}, adds to it, and writes it back. This
dependency means the loop cannot be pipelined to II\,=\,1 on its own
because the write of iteration $t$ must complete before iteration $t+1$
can read. The \texttt{UNROLL factor=4} directive resolves this by
unrolling four iterations into the same pipeline stage: four independent
partial products $x_j^2, x_{j+1}^2, x_{j+2}^2, x_{j+3}^2$ are computed
in parallel and their sum is added to \texttt{ss} once per group, reducing
the effective number of dependent additions by a factor of four.

Once \texttt{ss} has been computed, the scalar reciprocal RMS is available
and the second loop, \texttt{norm\_and\_scale}, has no loop-carried
dependencies. Every iteration reads one element of \texttt{x\_buff} and
one element of \texttt{weight\_buff}, multiplies by the scalar \texttt{ss},
and writes to \texttt{out\_buff}. These are entirely independent operations,
so the loop is amenable to wide unrolling. The \texttt{UNROLL factor=64}
directive allows up to 64 iterations to execute in parallel, subject to
memory port availability, which is satisfied by the \texttt{cyclic
factor=4} partitioning applied to all three local buffers.

\subsection{C-Simulation Results}
\label{subsec:rmsnorm_csim}

% Add csim pass/fail, testbench description, and numerical
% comparison against a software reference implementation.

\subsection{HLS Synthesis Results}
\label{subsec:rmsnorm_csyn}

% Add resource utilisation (LUT, FF, BRAM, DSP) and achieved II
% from the standalone rmsnorm synthesis report.

% =============================================================================
\section{Quantisation and Dequantisation}
\label{sec:quant_impl}
% =============================================================================

\subsection{Implementation}
\label{subsec:quant_impl}

% Describe quantize<S> and dequantize<S> from forward.h.
% Cover: group-based symmetric INT8 scheme, wmax scan, scale computation,
% round-and-clamp, why GS is used as a compile-time constant rather than
% runtime parameter, and the ARRAY_PARTITION pragmas.

\subsection{C-Simulation Results}
\label{subsec:quant_csim}

\subsection{HLS Synthesis Results}
\label{subsec:quant_csyn}

% =============================================================================
\section{Matrix--Vector Multiplication}
\label{sec:matmul_impl}
% =============================================================================

\subsection{Implementation}
\label{subsec:matmul_impl}

% Describe matmul<N,D> from forward.h.
% Cover: three-stage structure (load x/xs buffers, load w/ws per row,
% group accumulation), INT32 accumulation per group to prevent overflow,
% per-group dequantisation, ARRAY_PARTITION on x_buffer (factor=16) and
% w_buffer (factor=32), PIPELINE on outer row loop, commented-out UNROLL
% on inner loops and why they were left disabled in the prototype.

\subsection{C-Simulation Results}
\label{subsec:matmul_csim}

\subsection{HLS Synthesis Results}
\label{subsec:matmul_csyn}

% =============================================================================
\section{Softmax}
\label{sec:softmax_impl}
% =============================================================================

\subsection{Implementation}
\label{subsec:softmax_impl}

% Describe softmax<MAXSIZE> from forward.h.
% Cover: numerically stable max-subtract form, four sequential loops
% (max, exp, sum, norm), loop_tripcount hints for variable-length
% operation, UNROLL factor=16 on exp and norm loops, inv_sum optimisation.
% Note the template parameter change from <16> to <257> in forward_fyp.h
% and explain why (accommodates larger seq_len without code changes).

\subsection{C-Simulation Results}
\label{subsec:softmax_csim}

\subsection{HLS Synthesis Results}
\label{subsec:softmax_csyn}

% =============================================================================
\section{Grouped-Query Attention}
\label{sec:attn_impl}
% =============================================================================

\subsection{Baseline: Floating-Point Attention}
\label{subsec:attn_fp32}

% Describe the baseline attention in forward.cpp:
% - QKV projections via matmul
% - Float KV cache write via memcpy
% - Per-head score computation with PIPELINE on outer t-loop
%   and UNROLL on inner head_size dot product
% - softmax<16> call
% - Value accumulation with same outer-PIPELINE structure
% - Note that PIPELINE on outer t-loop is safe at seq_len=16 but
%   will cause deadlock at full scale (forward reference to Sec 4.7)

\subsection{Refined: Per-Head INT8 Attention}
\label{subsec:attn_int8}

% Describe the per-head quantisation approach in forward_fyp.cpp.
% Cover: outlier suppression motivation, three-step structure
% (quantise Q per head, load+quantise K on-the-fly, INT32 dot product,
% dequantise score), value accumulation with prob_int8 and v_int8,
% xb_accum with ARRAY_PARTITION complete (register file).
% Reference eq:int8_score and eq:int8_value_acc.

\subsection{C-Simulation Results}
\label{subsec:attn_csim}

\subsection{HLS Synthesis Results}
\label{subsec:attn_csyn}

% =============================================================================
\section{SwiGLU Feed-Forward Network}
\label{sec:ffn_impl}
% =============================================================================

% The FFN is commented out in the prototype kernel. This section will be
% written once the full TinyLlama kernel files (FFN_update2/update3) are
% incorporated. Content will cover:
% - Three-matrix SwiGLU structure (w1 gate, w3 up, w2 down)
% - Swish activation: SiLU approximation via sigmoid LUT (I-BERT strategy)
% - Element-wise gating and down projection
% - UNROLL factor=16 on gate loop
% - hb and hb2 buffer partitioning

\subsection{Implementation}
\label{subsec:ffn_impl_detail}

\subsection{C-Simulation Results}
\label{subsec:ffn_csim}

\subsection{HLS Synthesis Results}
\label{subsec:ffn_csyn}

% =============================================================================
\section{Top-Level Forward Function}
\label{sec:forward_impl}
% =============================================================================

\subsection{AXI Interface and Buffer Layout}
\label{subsec:forward_axi}

% Describe the AXI bundle configuration:
% - Prototype: 2 bundles (gmem0: transformer+key_cache+value_cache, gmem1: out)
% - Full model: 4 bundles (gmem0: weights, gmem1: key, gmem2: value, gmem3: out)
% Explain why sharing caches on gmem0 with weights is acceptable at small
% scale but causes deadlock at full scale.
% Cover static buffer declarations and ARRAY_PARTITION annotations.

\subsection{Layer Loop and KV Cache Management}
\label{subsec:forward_layers}

% Describe the main_forward_loop over L layers:
% - Embedding lookup (float table in prototype)
% - RMSNorm -> quantize -> QKV matmuls
% - KV cache write (float in prototype, INT8 in full model)
% - Multihead attention loop
% - Residual addition
% Note that FFN and final RMSNorm are commented out in the prototype.

\subsection{Known Limitations and Design Decisions}
\label{subsec:forward_limitations}

% Cover the three deliberate simplifications in the prototype:
% 1. Output projection W_O disabled (xb used directly in residual)
% 2. FFN commented out (attention-only forward pass)
% 3. Float KV cache (INT8 cache added in full model)
% Frame these as conscious scope decisions, not bugs.
